\section{Methods and Benchmark Design}
\label{sec:methods}

% Question: What do the methods of this project consist of?
% Answer: NOVA, the detector-level inference method, and the injector, which
% writes a known transit into the raw ramps of real JWST data so that
% recovery accuracy can be measured absolutely.
% Evidence: The structure of this section; NOVA_ASTRA_EXECUTION_20260905_R1/
% STATUS.md.

The methods have two parts. NOVA fits a forward model of the calibrated
detector time series directly to the retained pixels, with the transit
geometry inferred from the same data. The injector writes a known
transit into the raw ramps of a real out-of-transit JWST observation, so
that recovery accuracy can be measured against a known spectrum.
Appendix~\ref{app:nova} lists the constants of the
earlier release; where the present configuration differs, the value
given here is the current one.

\subsection{\nova{}}

% Question: What problem does NOVA solve, and in what form?
% Answer: It inverts a forward model of the detector: the same physical
% quantities that generate the pixels are inferred from them in one
% likelihood, on the real visit and on every injected input alike.
% Evidence: Sections 1.3 and 1.4; config/MATCHED_COVARIANCE_ENSEMBLE_R1.json
% (estimator.rebuild).

One forward model maps the transmission spectrum $D(\lambda)$, the limb
darkening $\bm{u}(\lambda)$, the stellar continuum, the spatial
response, and the additive background to every retained detector
sample. Spectrum, limb darkening, continuum, and background are inferred
from those samples in one likelihood, conditional on the response, the
orbit, and the source-curvature factor. Nothing is extracted, binned, or
fitted in a separate downstream stage. NOVA runs on the real WASP-17 visit, 720 integrations after native detector
processing, and on the injected 229-integration inputs of
Section~\ref{sec:injector}. Three calibration products are transferred
from the earlier release: the eight spatial background maps, the two
per-order uncertainty scales, and the form of the source-curvature
factor. Every other calibration product, and the white-light geometry,
is rebuilt on the input being analysed.

\subsubsection{The detector forward model}

% Question: What does NOVA predict for each detector sample?
% Answer: Continuum times exposure-integrated transit times spatial
% response, corrected by a source-curvature factor, plus an additive
% background.
% Evidence: proposals/NOMENCLATURE_PRESENTATION.md; proposals/
% handoff_nova_s_20260828/01_EXACT_NOVA_S_METHOD.md section 3.

Retained detector samples are organised into wavelength groups: one
group $g$ collects the pixels $P_g$ of one spectral order that share one
native detector column, and $g=g(p)$ denotes the group of pixel $p$. For
integration $t$ and pixel $p$, NOVA predicts
\begin{equation}
 \widehat Y_{tp}
 = \gamma_t\, C_{tg}\,
   \mathcal{T}_{tg}(D_g, \bm{u}_g;\, \Omega)\,
   q^\star_{pg}
 + B_{tp},
 \label{eq:nova-forward}
\end{equation}
and, where order supports overlap, the source term is summed over the
groups that contain the pixel. Here $C_{tg}$ is the continuum brightness
of the group at time $t$, and $\mathcal{T}_{tg}$ dims that light
according to the depth $D_g$ and limb darkening $\bm{u}_g$ at the
group's wavelength. The factor $\gamma_t$ corrects for slow source
variation over the visit, $q^\star_{pg}$ distributes the group's light
across its pixels with $\sum_{p\in P_g} q^\star_{pg}=1$, and $B_{tp}$ is
the additive background. Group geometry, wavelength assignment, and
response support follow the public trace and wavelength calibration
\citep{BainesEtAl2023Trace,BainesEtAl2023Wavelength} and the associated
spectral-response kernel \citep{DarveauBernierEtAl2022,JWSTPipeline2025}.
Because the background is never multiplied by the transit or by
$\gamma_t$, dimming it cannot imitate a change in depth. Valid samples have finite values,
positive uncertainties, and no \texttt{DO\_NOT\_USE} flag. On every
input, 97{,}557 pixels per integration are retained in 2{,}720 groups,
after six unsupported detector-edge groups are excluded.

\subsubsection{The shared spectrum and transit model}

% Question: How are the depths and the transit generated?
% Answer: One physical spectrum shared exactly by both orders, an exact
% Keplerian transit averaged over each exposure at the geometry from NOVA's
% own white-light fit, and small fitted limb-darkening deviations around a
% theory reference.
% Evidence: NOVA_SP_NOVA_S_SCIENCE_FREEZE_20260826_R1/configuration/
% variant_E12.json (outside the execution root); reports/
% ONEOVERF_CAUSAL_INVESTIGATION_20260909.md; NOVA_SP_NOVA_S_SCIENCE_FREEZE_
% 20260826_R1/code/project_sources/nonlinear_detector_sp_v0p1_20260720_src/
% nonlinear_sp/limb_darkening.py (lines 9-12, 935-965); reports/
% NOVA_LIMB_DARKENING_WORDING_CORRECTION_20260923.md.

Both orders evaluate the same spectrum on their own wavelength grids,
$D_g=D(\lambda_g)$, carried by $K=160$ adaptive basis coefficients,
\begin{equation}
  D_j = D_{\mathrm{global}}\,
  \frac{\exp\!\left[(\bm{H}_v\bm{v})_j\right]}
       {\sum_k w_k \exp\!\left[(\bm{H}_v\bm{v})_k\right]} ,
  \label{eq:depth-basis}
\end{equation}
which separates the achromatic depth $D_{\mathrm{global}}$ from a
weighted-zero-mean chromatic morphology and keeps all depths positive.
The order-discrepancy term $\Delta$ of the earlier release, an offset of
$\pm\Delta$ between the two orders' depths, is removed exactly; a bound
of $\pm10^{-12}$ on it breaks the invariance of the penalty derivatives.
Two depth slots without response support inside the validated wavelength
calibration are held as constants, so the nonlinear block holds 158
depth and eight limb-darkening coordinates. No atmospheric template,
wavelength-local smoother, or spectral prior is imposed.

The orbital state is
$\Omega=(P,\,a/R_\star,\,e,\,\omega_{\mathrm{peri}},\,b,\,t_0)$. Period
$P=3.73548546$~d, $e=0$, and $\omega_{\mathrm{peri}}=90^\circ$ are fixed
public values; $t_0$, $a/R_\star$, and the impact parameter $b$ come
from NOVA's own white-light fit on the same input
(Section~\ref{sec:own-white}) and are held fixed during the spectral fit.
An exact Keplerian engine \citep{Jaxoplanet2025} at backend quadrature
order 10 gives the instantaneous limb-darkened transit, and
$\mathcal{T}_{tg}$ is its average over the finite exposure with 21
Gauss--Legendre samples per integration. Quadratic limb darkening
\citep{MandelAgol2002} uses the coordinates of \citet{Kipping2013},
\begin{equation}
 u_1=2\sqrt{q_1}\,q_2, \qquad
 u_2=\sqrt{q_1}\,(1-2q_2).
 \label{eq:kipping-coordinates}
\end{equation}
Its theory reference is computed from STAGGER specific intensities
\citep{MagicEtAl2015} on the 6500~K plane, the closest complete plane to
$T_{\mathrm{eff}}=6550$~K, and interpolated bilinearly in $\log g$ and
$[\mathrm{Fe}/\mathrm{H}]$ to $\log g=4.149$ and
$[\mathrm{Fe}/\mathrm{H}]=-0.25$. The reference file is pinned by
content hash. Per coordinate and
order the fit has one offset and one normalised log-wavelength slope,
eight coordinates in all, penalised for deviation from theory with scale
0.5 in logit space and for curvature with scale
$2.0\,\mu\mathrm{m}^{-2}$. Fine nodes outside the theory support are
excluded, not extrapolated. Only the spectral fit uses the theory
reference; the white-light fit leaves limb darkening free.

\subsubsection{The empirical spatial response}

% Question: Where does q-star come from?
% Answer: The transit itself reveals where the starlight lands: a per-pixel
% Huber regression on a depth-free transit template, with amplitudes
% normalised within each group, rebuilt on every input.
% Evidence: products/original_v45_reference_qstar_r1_package/source/
% build_v47_empirical_template.py and aggregate_v47_qstar.py; products/
% original_A_group_calibration_adaptation_r1/build_nova_s_real_wasp17_responses_r1.py.

The spatial response is estimated from the data themselves. Starlight
dims during transit and the additive background does not, so the
transit-correlated amplitude of each pixel measures how much starlight
lands there. For each order, a depth-free
event template $\phi_o(t)$ is the median of baseline-normalised
bright-pixel residuals over all transit phases. Every qualifying pixel is
then fitted with Huber weights as
\begin{equation}
 Y_{tp} = A_{pg}\,\phi_{o}(t) + m_p + n_p\,\tau_t + \varepsilon_{tp},
 \label{eq:qstar-fit}
\end{equation}
with a level $m_p$, a drift $n_p\tau_t$, and the transit-correlated
amplitude $A_{pg}$. Pixels qualify when at least 95\% of their
out-of-transit samples are valid, their baseline is positive, and they
lie above the 35th percentile of baseline brightness. The response is the
normalised amplitude,
$q^\star_{pg}=A_{pg}/\sum_{p'\in P_g}A_{p'g}$, so total group brightness
remains with the continuum. Two alternating folds, stratified across the
transit phases, provide independent estimates. A second-difference
smoother along detector column and integer trace rank takes its strength
from the opposite fold, and the response is accepted only if the folds
agree in flux distribution and centroid (Appendix~\ref{app:nova}). Its
integer trace-relative smoothing coordinate produces occasional centroid
resets and an amplitude excess at the reddest columns, which biases the
recovered transit shallow there.

\subsubsection{The continuum and the additive background}

% Question: How are the stellar continuum and the background modelled and
% calibrated?
% Answer: A level and slope per group for the continuum; eight fixed spatial
% maps crossed with two temporal patterns for the background, the temporal
% pattern and anchor recomputed on the off-trace pixels of each input.
% Evidence: build_e13b_basis_products.py (cluster-only); products/
% original_A_group_calibration_adaptation_r1/build_nova_s_real_wasp17_e13b_r1.py;
% reports/OFFTRACE_TRANSIT_COMPONENT_DIAGNOSIS_20260921.md section 2.

Each group's continuum is linear in time,
$C_{tg}=\beta_{0g}+\beta_{1g}\tau_t$, giving 5{,}440 linear coefficients.
The background is a sum of eight fixed spatial maps $\Phi_{pk}$ with two
temporal patterns,
\begin{equation}
 B_{tp}=\sum_{k=1}^{8}\Phi_{pk}
 \left[\alpha_{k0}+\alpha_{k1}\,G_1(t)\right],
 \label{eq:background-model}
\end{equation}
sixteen coefficients in all. The maps, a standardised calibration
background anchor and seven low-order polynomial surfaces orthonormalised
on off-trace pixels, were chosen by blocked off-trace cross-validation
on the earlier release and are transferred unchanged. Off-trace pixels
of each input supply the rest: per-integration coefficients, the
temporal pattern $G_1(t)$, and an expectation
$\bm{\alpha}_{\mathrm{off}}$ with precision $\bm{N}_{\mathrm{off}}$.
$G_1$ is the first of the eight leading principal components of the
off-trace time series whose absolute correlation with the empirical
transit template is at most 0.25. The
expectation enters the objective as a soft anchor, and the background is
refitted jointly with the continuum at every step. No fixed background is subtracted at any stage, the likelihood
contains no Gaussian-process term, and outliers are handled by the Huber
weights below.

\subsubsection{The source-curvature factor and detector uncertainties}

% Question: What are gamma_t and sigma_tp?
% Answer: An out-of-transit curvature correction of fixed common form,
% recomputed on each input's alternating folds, and per-sample uncertainties
% scaled per order and per detector row from out-of-transit cross-validation.
% Evidence: products/original_A_group_calibration_adaptation_r1/
% build_nova_s_real_wasp17_source_curvature_r1.py and
% build_nova_s_real_wasp17_response_oot_r1.py; config/
% CURVATURE_POLICY_AUTHORIZATION_20260919.json; NOVA_SP_NOVA_S_SCIENCE_FREEZE_
% 20260826_R1/code/runtime/N_SCALE_CONTRACT.json (outside the execution root);
% reports/offtrace_transit_diagnosis_20260921/sources/v47_oot_guard.py.

The factor $\gamma_t$ corrects for slow variation of the source over the
visit. Stable trace flux is integrated per order over the out-of-transit
integrations, and the ratio of a weighted quadratic in time to its affine
approximation defines the factor, normalised to unit out-of-transit mean.
Its common form was selected on the earlier release by held-out
chi-squared on two folds against null and per-order alternatives
(Appendix~\ref{app:nova}) and is transferred. On each input the factor is
recomputed from that input's out-of-transit integrations in alternating
folds (175 and 175 for the real visit). It is accepted only if the fold
event-mean difference is at most 250~ppm and the full excursion at most
1000~ppm. A fold correlation of at least 0.95 is a gate for the real
reference and a reported diagnostic for ensemble realisations. The
factor multiplies only the stellar term of Equation~\ref{eq:nova-forward}
and adds no fitted parameters.

Per-sample uncertainties are
$\sigma_{tp}=\mathrm{ERR}_{tp}\,N_o\,s_{r(p)}$: the calibrated ERR array,
scaled once per order by the $N_o$ transferred from the earlier release
and once per detector row by $s_{r(p)}\geq 1$. The row scale comes from
bidirectional out-of-transit cross-validation of a level and slope per
row, rebuilt on each input (280 calibration and 70 held-out integrations
on the real visit).

\subsubsection{The complete objective}

% Question: What does NOVA minimise?
% Answer: A Huber detector mismatch over the retained mask plus the
% off-trace background anchor and the limb-darkening penalties.
% Evidence: NOVA_SP_NOVA_S_SCIENCE_FREEZE_20260826_R1/code/sparse_varpro/
% production_irls.py and code/project_sources/nonlinear_detector_sp_v0p3_
% huber_discrepancy_factorial_20260722_r4_src/nonlinear_sp_huber_factorial/
% robust.py (HUBER_DELTA 1.345), both outside the execution root.

With the standardised residual
\begin{equation}
 r_{tp}=\frac{Y_{tp}-\widehat Y_{tp}
 (D,\bm{u},\bm{\beta},\bm{\alpha};\,q^\star,\Omega)}{\sigma_{tp}},
 \label{eq:residual}
\end{equation}
in which fitted parameters appear before the semicolon and fixed inputs
after it, the estimate is
\begin{equation}
\begin{split}
 &(\widehat D,\widehat{\bm{u}},\widehat{\bm{\beta}},
  \widehat{\bm{\alpha}})
 = \arg\min\Big\{ \\
 &\quad\sum_{(t,p)\in\mathcal{M}}\rho_{1.345}(r_{tp}) \\
 &\quad+\tfrac{1}{2}(\bm{\alpha}-\bm{\alpha}_{\mathrm{off}})^{\mathsf T}
   \bm{N}_{\mathrm{off}}(\bm{\alpha}-\bm{\alpha}_{\mathrm{off}}) \\
 &\quad+\tfrac{1}{2}\|\bm{r}_{\mathrm{LD,theory}}(\bm{u})\|_2^2
  +\tfrac{1}{2}\|\bm{r}_{\mathrm{LD,curv}}(\bm{u})\|_2^2
 \Big\},
\end{split}
 \label{eq:objective}
\end{equation}
where $\mathcal{M}$ is the retained sample mask and $\rho_{1.345}$ the
Huber loss. The three terms are the detector mismatch, the off-trace
background anchor, and the limb-darkening theory and smoothness
penalties.

\subsubsection{The nested solver}

% Question: How is the objective minimised, and what counts as convergence?
% Answer: Huber weights outside, bounded trust-region steps in the middle,
% an exact linear solve inside; acceptance requires convergence gates and
% agreement of independent starts.
% Evidence: NOVA_SP_NOVA_S_SCIENCE_FREEZE_20260826_R1/code/sparse_varpro/
% production_irls.py (outside the execution root); reports/
% PHYSICAL_ENDPOINT_THIRD_START_20260918.md; reports/
% LOW_UNCAL_V7_SPECTRAL_CLOSURE_ARCHIVE_20260923.md.

The parameters split into a nonlinear block $\bm{\theta}=(D,\bm{u})$ and
a linear block $(\bm{\beta},\bm{\alpha})$, and the solver nests three
layers. Innermost, at fixed $\bm{\theta}$ and weights, variable
projection \citep{GolubPereyra1973} solves the linear problem exactly,
\begin{equation}
\begin{split}
 (\widehat{\bm{\beta}},\widehat{\bm{\alpha}})=
 \arg\min_{\bm{\beta},\bm{\alpha}}\;
 &\frac{1}{2}\sum_{(t,p)\in\mathcal{M}}\frac{w_{tp}}{\sigma_{tp}^2}
 \left[Y_{tp}-\widehat Y_{tp}\right]^2 \\
 &+\frac{1}{2}(\bm{\alpha}-\bm{\alpha}_{\mathrm{off}})^{\mathsf T}
   \bm{N}_{\mathrm{off}}(\bm{\alpha}-\bm{\alpha}_{\mathrm{off}}),
\end{split}
 \label{eq:varpro}
\end{equation}
by factorising the sparse continuum block and reducing the background
coupling to a sixteen-dimensional Schur complement. The profiled
residual and Jacobian differentiate through this exact inner optimum
(Appendix~\ref{app:nova}). In the middle, bounded trust-region reflective
least squares \citep{BranchColemanLi1999} minimises the profiled
objective over $\bm{\theta}$. Outermost, Huber reweighting
\citep{Huber1964} updates $w_{tp}=\min(1,\,1.345/|e_{tp}|)$ for
standardised residual $e_{tp}$, holding the weights fixed during each
trust-region solve.

A fit is accepted only when the weights and the objective have stopped
changing and the shared spectrum is stationary to within 0.01 and
0.05~ppm on the two spectral summaries. First-order optimality must be
met, a confirmation fit at the recomputed final weights must reproduce
the solution, and a descent audit must find no material improvement
along scaled Cauchy and Gauss--Newton directions
(Appendix~\ref{app:nova}). Each analysis starts at least twice,
independently: at the white-light depth with zero morphology and the
theory limb darkening, and with a fixed 250~ppm morphology perturbation
and small limb-darkening offsets. These two starts must agree to within
0.1 and 0.5~ppm on the same summaries; injected recoveries add a third
start at the midpoint of the two start vectors. Among the starts that
pass every gate, the one with the lowest Huber objective is selected.

\subsubsection{White-light geometry inference}
\label{sec:own-white}

% Question: Where do t0, a/R* and b come from?
% Answer: From NOVA's own joint two-order white-light fit on curves built
% from its own detector data, sampled with nested sampling through public
% Eureka!, accepted by a grazing-mass rule fixed before calibration, and
% handed over as whole-posterior weighted medians.
% Evidence: source/nova_own_white_eureka.py, source/
% nova_own_white_nested_20260922_r5.py, source/nova_own_white_geometry.py,
% source/v7_grazing_statistic_20260923.py; config/
% V7_GEOMETRY_GATE_REGISTRATION_20260923.json; reports/V7_CALIBRATION_RESULT_20260923.md;
% products/LOW_white_overlay_20260921_r1/uncal/own_white.epf.

NOVA infers $t_0$, $a/R_\star$, and the inclination $i$, hence
$b=(a/R_\star)\cos i$ for $e=0$, from its own detector data before the
spectral fit, without any value from another pipeline or from the
injector. For each order a white-light curve is built per integration
from the retained detector samples of that order on the same input. The
two curves are fitted jointly with the public Eureka! light-curve stage
\citep{BellEtAl2022} as two channels sharing one geometry. Common to
both are $t_0$, $a/R_\star$, and $i$. Free per order are the radius
ratio $k_o$, two quadratic limb-darkening coefficients
\citep{MandelAgol2002}, a linear baseline, and a white-noise multiplier,
with $P$, $e$, and $\omega_{\mathrm{peri}}$ fixed as above. Priors are
wide and independent of any fit. $t_0$ is uniform over the observed
window, $a/R_\star$ uniform on $[1.01,\,100]$, $i$ uniform on
$[0^\circ,\,90^\circ]$, and $k_o$ normal with mean 0.1255 and width
0.085. Limb-darkening coefficients are uniform on $[0,1]$, baseline
level and slope normal about 1 and 0 with widths 0.05 and 0.01, and the
noise multiplier normal with mean 1.4 and width 2.0. The posterior is
sampled with dynesty \citep{Speagle2020} through Eureka!: 2{,}048 live
points, multi-ellipsoid bounds, random-walk proposals, and a stop when
the remaining evidence $\Delta\ln Z$ reaches 0.01. Nested sampling is
used throughout because the affine-invariant ensemble sampler
\citep{ForemanMackeyEtAl2013} did not reach the required 50
autocorrelation lengths on the raw-level inputs.

Acceptance follows a rule fixed before it was calibrated. For every
posterior sample, $b$ is computed from that sample's $a/R_\star$ and
$i$, and the sample is grazing or non-transiting if
\begin{equation}
 b > 1-\max(k_1,\,k_2).
 \label{eq:grazing}
\end{equation}
Every sample and its native weight are kept. The grazing mass is the
weighted fraction of such samples, with a central 95\% interval from 200
native resampling replicates of the run and a maximum half-width of
0.01. When no sample of positive weight is grazing, the upper bound is
$\min(1,\,3/\lfloor N_{\mathrm{ESS}}\rfloor)$ with $N_{\mathrm{ESS}}$ the
Kish effective sample size, an approximation for dependent nested
samples. A fit passes if the upper bound is below 1\%, and passes with
reported grazing mass between 1\% and 5\%. It is held at 5\% or more,
if the sampler did not reach its stop, or if the interval is too wide.
Held fits are not continued, reseeded, or trimmed, and no geometry is
imported in their place. On acceptance the handoff is the weighted
median of the whole posterior for $t_0$, $a/R_\star$, $i$, and each
$k_o$; the white depth per order is $10^6 k_o^2$, and the spectral fit
starts at their equal-order mean. Multimodality and tail diagnostics are
reported and do not gate. Before use, the rule was calibrated on saved
posteriors: chains that had failed their length assessment must be held,
and accepted chains and the real nested posterior must pass.

\subsubsection{The uncertainty ensemble}
\label{sec:ensemble}

% Question: How are NOVA's spectral uncertainties obtained and validated?
% Answer: By refitting from scratch on realisations built from the accepted
% real model plus resampled real residual blocks; the ensemble errors define
% a finite-ensemble Student likelihood that must pass held-out coverage.
% Evidence: config/MATCHED_COVARIANCE_ENSEMBLE_R1.json; products/
% generator_v2_20260923_v7_r2_package (REGISTRY.json, source/
% generator_v2_draw_20260922.py, pilot_noise.py); source/
% finite_ensemble_likelihood.py; config/COVARIANCE_ANALYSIS_PROTOCOL_R1.json;
% reports/COVARIANCE_VALIDATION_RESULT_20260921.md.

The profiled normal system gives a covariance conditional on one fixed
set of calibration products. NOVA therefore estimates its spectral
uncertainty by repeating the whole analysis on realisations of the real
visit,
\begin{equation}
 Y^{(k)}_{tp}=\widehat Y^{\mathrm{ref}}_{tp}+\bar r_p+a_p\,\psi_t
 +\mathrm{ERR}_{tp}\,z_{d_k(t),\,p},
 \label{eq:ensemble-draw}
\end{equation}
where $\widehat Y^{\mathrm{ref}}$ is the accepted final detector
prediction of the real visit on the trace pixels and the fitted
background mean off the trace. The term $\bar r_p$ is the uncentered
per-pixel mean of the real out-of-transit residual. The term $a_p\psi_t$
projects the amplitude map of the vetoed leading off-trace temporal
component onto the transit decrement of the accepted white-light fit,
$\psi_t=\mathcal{T}^{\mathrm{w}}_t-1$, which is exactly zero out of
transit. The last term is the real ERR array times a centred
standardised out-of-transit residual $z$ from donor integration
$d_k(t)$. Donors are whole frames drawn in circular blocks of 64
integrations within one of the two out-of-transit intervals, $[0,230)$
and $[600,720)$, never across the transit. The same schedule feeds the
trace and off-trace pixels, and the real ERR and quality arrays are
kept.

The ensemble has 352 realisations with fixed seeds and roles: 256 for
estimation, 32 held out for validation, and 32 each at block lengths 128
and 192 with the held-out seeds. Each is refitted from scratch: fresh
spatial response, background pattern and anchor, curvature factor and
row guard, its own nested white-light fit under the rule of
Section~\ref{sec:own-white}, then the two strict starts. A realisation
whose white fit is held or whose starts are ineligible is recorded as a
failure and stays in the denominator.

A realisation's error is its recovered spectrum minus the reference
spectrum in the $p=128$ measured modes of the reporting operator. From
the $N=256$ estimation errors the sample mean $\bm{m}$ and unbiased
covariance $\bm{S}$ are formed without shrinkage, jitter, or mode
selection. The likelihood of a new error $\bm{x}$ is the finite-ensemble
multivariate Student distribution \citep{SellentinHeavens2016} with
$N-p$ degrees of freedom and scale
\begin{equation}
 \bm{\Sigma}=\left(1+\frac{1}{N}\right)\frac{N-1}{N-p}\,\bm{S},
 \label{eq:finite-ensemble}
\end{equation}
centred on $\bm{m}$, so that
$(\bm{x}-\bm{m})^{\mathsf T}\bm{\Sigma}^{-1}(\bm{x}-\bm{m})/p$ follows
$F(p,\,N-p)$. Validation uses the 32 held-out realisations. The
likelihood is flagged if joint 95\% coverage is rejected by a one-sided
exact binomial test at $p<0.025$, or if Mardia skewness or kurtosis
falls outside 512 Gaussian reference simulations at $p<0.025$.
Validation runs once all 352 realisations have finished, and the
retrieval likelihood follows only if validation passes.

\subsection{The injector}
\label{sec:injector}

% Question: How is a truth-known benchmark built from real JWST data?
% Answer: A known transit signal is written into the raw group ramps of real
% out-of-transit integrations, dimming only modelled target starlight, so
% that detector processing, background and noise are real and the truth is
% known.
% Evidence: source/emit_original_A_group_segment.py, source/
% original_v45_group_signal.py, source/group_accumulated_transport.py;
% evidence/matrix_trip_bindings_r2.json.

The injector answers the validation problem of
Section~\ref{sec:detector-to-validation}: real observations provide no
ground truth, so a truth is hidden inside real data. It works on raw
detector ramps rather than on calibrated integrations, so that every
recovery method runs its own detector processing, including group-level
$1/f$ correction, on the same injected frames.

\subsubsection{Baseline and null}

% Question: Which real data carry the injected signal, and what is the
% control?
% Answer: The first 229 raw integrations of the WASP-17 visit, all before
% the real transit; the null is the same frames through the same processing
% with nothing injected, and a rate-level control adds the signal after
% ramp fitting.
% Evidence: evidence/matrix_trip_bindings_r2.json (baseline_provenance);
% reports/PHYSICAL_MATCHED_NATIVE_NULL_20260918.md; products/
% w17_measured_raw_emission_20260923_v7_r3_prepared/source/
% emit_conditional_matched_rate_20260920.py.

The baseline is the raw, group-level NIRISS/SOSS observation of WASP-17
(program 1353, observation 101, visit 1), restricted to its first 229
integrations, in two raw segments of 178 and 51. Because the real
transit and its buffers occupy integrations 230 to 599 of the
720-integration visit, every baseline integration is pre-transit data
with no event present. Each integration holds eight non-destructive
reads of the $256\times2048$ subarray at a group time of 5.494~s, one
frame per group and no dropped frames, so the accumulated exposure times
run from 5.494 to 43.952~s. The null is the same 229 raw integrations,
with nothing injected, through the identical detector chain. A
rate-level control adds the signal after ramp fitting instead. There,
the count change of the final read, divided by its exposure time and
projected without the flat-field reversal, is added to the calibrated
rate images of the null. Their uncertainty and quality arrays are
untouched. The difference between raw-level and rate-level recoveries
isolates what group-level processing does to the injected signal.

\subsubsection{What light is dimmed}

% Question: Which photons transit?
% Answer: A map of modelled target starlight in three zones by distance from
% the traces: a fixed core, measured wings, and a far region injected at a
% fraction of the measured residual light, with two bracketing assumptions
% and a measured per-column member.
% Evidence: products/light_based_family_20260921_r1/FAMILY.json and
% FAR_PIXEL_ACCOUNTING.json; products/full_light_outer_shape_controls_20260916_r1/
% SOURCE_ESTIMATE.json; products/measured_per_column_20260922_r1/REGISTRATION.json;
% source/freeze_measured_columns_20260922.py; source/
% analyze_evening_rate_far_fraction_20260921.py; products/
% w17_measured_raw_emission_20260923_v7_r3_prepared/source/
% fixed_primary_family_signal_20260920.py; reports/ENDPOINT_HOLD_VERIFICATION_20260920.md.

Only modelled target starlight transits. The source is a map in
$\mathrm{DN\,s^{-1}}$ at a reference time, in three zones by the
distance $d$ from a pixel to the nearer saved trace. The core,
$d\leq16$ pixels (119{,}277 pixels), is an earlier projection of the full
target photon budget with its saved calibration factors, kept byte for
byte. The wings, $16<d\leq80$ pixels (257{,}554 pixels), are the
background-subtracted out-of-transit light measured on the same data
without smoothing, per 64-column band and row parity, under a
nonnegativity constraint that preserves each band's signed sum. The far
region, $d>80$ pixels (121{,}811 pixels), holds the residual light after
background subtraction. Whether that light belongs to the target is not
established, so it is injected at a fraction $m$ of the measured field.
Two bracketing assumptions define the family, none of the far light
($m=0$) and all of it ($m=1$), with a midpoint member ($m=0.5$) used only
in the delivery check. Far pixels with zero measured light (33{,}758)
receive zero.

A fourth member uses the measured far-light fraction. On the final
calibrated images of the full real visit, the mean signal of the
102{,}308 safe far pixels drops during the real transit by 0.834\% over
all columns. The measurement uses 350 out-of-transit and 315 in-transit
integrations and excludes 55 buffer integrations. Per detector column, the dip is
divided by the real white transit depth of the same integrations as
accepted when the member was fixed, 0.0136288893. The currently accepted
white fit gives 0.0136285918, a difference of $3\times10^{-7}$ that
leaves the member unchanged. The ratio is averaged over a 128-column
window using only the 1{,}884 columns with a direct measurement and
bounded to $[0,1]$. The resulting multiplier $m_x$ has minimum 0.103,
median 0.606, and maximum 1, with 84 columns at the upper bound, and
carries 61.5\% of the full far-light field. Because it is measured
after native $1/f$ correction, which may imprint the far-row statistic,
it is an upper bound on the stellar far fraction, not a measurement of
it. Pixels flagged by the F277W exposure (12{,}808) keep the modelled
target light and are excluded from the empirical estimates; the
10{,}208 reference pixels receive none.

The core keeps its earlier wavelength and order kernels. The outer zones
take the wavelength and order mixture of a separately bound outer model
through a per-pixel factor. The outer light thus has a spectrum and a
time derivative but not the core's brightness. Both components are
projected separately and summed, and the sum must reproduce the fixed
map to a relative $5\times10^{-13}$ (absolute
$10^{-10}\,\mathrm{DN\,s^{-1}}$). In the core, the stellar spectrum
$S(\lambda,\tau)$ varies linearly in time between two calibrated
endpoint spectra and is held at the endpoint value outside its
calibration interval, which affects only the reads of integration 0. The
outer component is affine between endpoints defined on the full physical
read interval.

\subsubsection{The per-read transit signal}

% Question: How is the transit computed?
% Answer: With the ExoTETHyS TRIP kernel on MPS-ATLAS intensities, averaged
% by Gauss-Legendre quadrature over each of the eight reset-to-read windows
% of every integration, at the orbit of the public control file.
% Evidence: products/original_A_causal_validation_trip_reference_r1_package/
% source/run_original_v45_group_trip.py (per-read packets); source/
% B_mixed_native_trip.py (integration-level packets, 5e-13 edge gate);
% products/w17_measured_raw_emission_20260923_v7_r3_prepared/source/
% approved_endpoint_hold_20260920.py; source/original_v45_group_signal.py
% (group_windows); evidence/matrix_trip_bindings_r2.json (geometry);
% config/CURVATURE_PARTITION_CONTROL_R1.json.

The injected spectrum $D(\lambda)$ is defined on the injector's own
15{,}326-node wavelength grid, and the radius ratio at each node is
$\sqrt{D(\lambda)}$. The occulted flux comes from the TRIP routine of
ExoTETHyS \citep{MorelloEtAl2020} over 20{,}000 annuli, under a
regression test against the public implementation. It uses
limb-darkened intensities from the MPS-ATLAS grid (set 1)
\citep{KostogryzEtAl2022,KostogryzEtAl2023} for $T_{\mathrm{eff}}=6550$~K,
$\log g=4.149$, and $[\mathrm{M}/\mathrm{H}]=-0.25$. The injected orbit
is the circular orbit of the public control file of the published
WASP-17~b analysis \citep{LouieEtAl2025}, with the transit duration
matched to that file; no recovery receives these values. Within
integration $t$, read $n$ closes a window from the reset to the end of
that read, and the transit is averaged over each of the eight nested
windows by Gauss--Legendre quadrature. Grids of 41 and 61 nodes must
agree to better than 1~ppm, and the 61-node result is used. Two boundary
conditions apply. In the integration-level packets of the rate-level
control, the first and last integrations must be unocculted to
$5\times10^{-13}$. In the per-read packets, integration 0, the only one
whose read windows cross the calibration interval, must carry an exactly
zero time moment. Its remaining fraction, a null round-off
$2\times10^{-16}$ below unity, is set to exactly 1. Per integration,
read, and node, the packet stores the window mean
$\langle T\rangle_{tn}(\lambda)$ and the first time moment
$\mu_{tn}(\lambda)=\langle(\tau-\tau_t)(T-1)\rangle_{tn}$ about the
integration midpoint $\tau_t$. On the 229-integration timeline the
atmospheric packets carry a nonzero decrement on integrations 50 to 178
(129 integrations, 100 out of transit). The grey packet of the delivery
check uses the narrower window 51 to 177 (127 integrations).

\subsubsection{Transport to raw counts}

% Question: How does the signal enter the raw ramps?
% Answer: As an exact accumulated count change per read, projected to the
% detector, returned to unflattened counts, and written through the inverse
% of the native linearity correction.
% Evidence: source/original_v45_group_signal.py; source/
% group_accumulated_transport.py; config/DEVELOPMENT_GROUP_TRANSPORT_R1.json;
% evidence/original_group_progress_r3.json; evidence/
% exo250_raw_configuration_bindings_20260923_r1.json (step sequence, scale 1.0,
% flat range); stage3b_group_signed_saturation_adapter_r2.py (cluster-only).

For a source linear in time, the count change accumulated by read $n$ of
integration $t$ at node $\lambda$ is the exact integral of $S\,(T-1)$
over the window,
\begin{equation}
\begin{split}
 \Delta N_{tn}(\lambda) = {}& s_n\big[S_t(\lambda)\,
 \big(\langle T\rangle_{tn}(\lambda)-1\big) \\
 &\quad +\dot S(\lambda)\,\mu_{tn}(\lambda)\big], \\
 \Delta Y^{\mathrm{raw}}_{tnp} = {}& f_p
 \sum_{o}\big[\bm{A}_o\,\Delta N_{tn}\big]_p .
\end{split}
 \label{eq:v45-injector}
\end{equation}
Here $s_n=5.494\,n$~s is the accumulated exposure time of read $n$, the
same for every integration, $S_t$ the stellar spectrum at the
integration midpoint, and $\dot S$ its slope. $\bm{A}_o$ is the detector
projection of order $o$: the sparse spectral-response matrix from the
public kernel, a wing correction, a per-column normalisation, and a
per-pixel out-of-transit response calibrated on the data. $f_p$ is the
flat-field coefficient of the bound native reference file,
\texttt{jwst\_niriss\_flat\_0275} (0.2021 to 1.5027 on the original
507{,}606-pixel two-order support, quality flags left in place), which
returns the flat-fielded loss to raw counts. The change is zero on the
reference pixels and must be non-positive everywhere, or the emission
stops. It is written into the ramp in raw DN. An affine bridge carries
the raw counts to the input of the native linearity step, after the
pipeline's group-level preprocessing up to the reference-pixel
correction. A Newton solve of the pipeline's linearity polynomial then
finds the raw count whose linearised value is the original one plus
$\Delta Y^{\mathrm{raw}}$, and the inverse bridge returns to the raw
file. No gain conversion is applied. The round trip must close to
0.015625~DN, the inverse linearity to $10^{-8}$, and the per-order
accounting exactly. Samples flagged saturated in the native group
quality array, or non-finite at the linearity input, keep their values,
and the undelivered decrement is recorded. Observed noise, background,
and field sources are untouched, and no scene mask enters the emission.

\subsubsection{Native detector processing}

% Question: What happens to the injected ramps before recovery?
% Answer: NOVA's own input runs the public group-level background and 1/f
% steps with the F277W-derived mask, then the calibration pipeline's
% linearity, jump, ramp-fitting and gain steps, then Stage 2; in production
% each comparator runs its own Stage 1 on the same raw frames.
% Evidence: source/run_nova_original_A_group_detector.py; reports/
% PHYSICAL_MATCHED_NATIVE_NULL_20260918.md; evidence/
% native_null_completion_20260918_r1.json (NATIVE_F277W_MASK_USED.json is
% cluster-only).

NOVA's own input is prepared from the injected raw segments by the public
chain it has always used. The initial group steps of the calibration
pipeline \citep{JWSTPipeline2025} run through reference-pixel
correction, followed by the exoTEDRF group-level background and $1/f$
steps \citep{Radica2024}. The pipeline's linearity, jump, ramp-fitting,
and gain steps and Stage 2 with the flat field then give rate images in
$\mathrm{DN\,s^{-1}}$. The $1/f$ step masks the contaminating sources
identified in the F277W exposure, and the mask built at run time must
equal the calibrated control of 12{,}808 pixels. No project scene mask
is applied, and the null follows the same chain. Pilot comparisons are
rate-level: each comparison pipeline enters at its native extraction
after this common Stage 2. In production each will consume the same
injected raw frames and run its own Stage 1 with its own group-level
$1/f$ treatment.

\subsubsection{Delivery checks}

% Question: How is it known that the requested signal reached the data?
% Answer: Read-count, positivity and protected-pixel closures, exact native
% aperture sums, and a grey linearity test: a constant 16,000 ppm depth
% injected under none, half and all of the far light must respond linearly
% to within 10 ppm equivalent depth per column.
% Evidence: reports/LIGHT_BASED_GREY_DELIVERY_RESULT_20260921.md;
% reports/MEASURED_GREY_DELIVERY_RESULT_20260922.md; products/
% light_based_family_20260921_r1/FAMILY.json (delivery).

Every emission must close its read-count, nonnegativity, and
protected-pixel checks and reproduce the native aperture sums of the
processed frames exactly. The family is then tested with a grey signal:
a constant depth of 16{,}000~ppm at every node, injected under the three
far-light assumptions and processed natively. Per detector column, the
aperture response at $m=0.5$ must lie within 10~ppm equivalent depth of
the secant between the responses at $m=0$ and $m=1$, over 2{,}188
identified columns. Delivery passes only if the maximum departure is
below 10~ppm. The measured-fraction member is checked against the
prediction of the column-local secant at its own $m_x$ under the same
bound. A pass shows that the delivered response is consistent across the
family; it does not show that the far light is stellar.

\subsubsection{The endpoint protocol and envelope rule}

% Question: How is the far-light uncertainty carried into a result?
% Answer: The reference atmosphere is recovered at both bracketing
% assumptions; the spread between its two raw-level recoveries decides the
% wording of the family's results, not their release.
% Evidence: products/light_based_family_20260921_r1/FAMILY.json
% (endpoint_criterion, chromatic.cases); reports/
% PHYSICAL_INJECTOR_PROTOCOL_20260918.md; reports/
% OVERNIGHT_PRODUCTION_ADDENDUM_20260921.md; reports/
% W17_PRODUCTION_BINDING_INVENTORY_20260923.md.

The reference atmosphere is recovered at both bracketing assumptions, at
the raw level and at the rate level. Its envelope is the root-mean-square
difference between the two raw-level recoveries over the 147 reporting
bins. At or below 10~ppm a result of the family is reported as validated
within the scope of the source family; above 10~ppm it is reported as
usable with its envelope. The envelope changes the wording, not the
release. The two endpoints bracket the far-light assumption; they are
not confidence bounds. The 10~ppm envelope is set against the 30~ppm
difference between pipelines that the comparison is designed to resolve.
A second pilot atmosphere is recovered at the full-far-light endpoint
only, at the raw level. Each pilot case has fresh calibration products,
its own white-light geometry, and three strict starts. Production
atmospheres are recovered once each at the measured-fraction member and
carry the pilot envelope wording: 24 atmospheres per visit, twelve with
local condensation chemistry and twelve with rainout chemistry, with two
strict starts each. Six visits are in the programme, of WASP-17,
HAT-P-26, WASP-39, HAT-P-18, WASP-96, and WASP-80. Each is released only
after its own accepted geometry, core and wing checks, and grey delivery
under the same rule.

\subsubsection{Comparison analyses and fairness}

% Question: Against what is NOVA compared, and on what terms?
% Answer: Three public pipelines run end to end as their authors run them,
% on the same injected frames, each with its own white-light inference;
% only the three geometry prior lines are shared.
% Evidence: reports/COMPARATOR_AUTHOR_PROTOCOL_AUTHORIZATION_20260923.md;
% config/COMPARATOR_SHARED_GEOMETRY_PRIORS_20260923_R1.json; reports/
% TRANSITSPECTROSCOPY_REMAINING_CONFIGURATION_20260922.md; reports/
% TS_MPS2_SPECTRAL_CONTINUATION_PREPARATION_20260923_R1.md.

Three public pipelines run as native analyses on the same injected
frames: the Ahsoka analysis of the published WASP-17~b spectrum
\citep{LouieEtAl2025}, exoTEDRF 2.5.0 \citep{Radica2024}, and
transitspectroscopy 0.3.12 \citep{Espinoza2022TransitSpectroscopy}. Each
keeps its authors' sampler, run length, fit structure, and acceptance.
Ahsoka fits the two orders separately with the Eureka! ensemble sampler
\citep{ForemanMackeyEtAl2013} at 200 walkers, 1{,}000 steps, and 500
burn-in. It fixes each order's native median $t_0$, $a/R_\star$, and $i$
into that order's spectral fit. exoTEDRF fits a joint two-order white
curve with its native nested sampler and 13 free parameters.
transitspectroscopy traces the median image with its double-Gaussian
template and extracts a box aperture of 15-pixel radius. Its joint
two-order white curve is fitted with juliet \citep{EspinozaEtAl2019}
nested sampling, 17 free parameters and a Mat\'ern-3/2 process per
order; its spectral fit uses ExoTETHyS limb darkening from MPS-ATLAS set
2 in Kipping coordinates. In every white fit only three prior lines are
replaced, by the priors NOVA uses for this observation. These are $t_0$
uniform over the observed window, $a/R_\star$ uniform on $[1.01,\,100]$,
and $i$ uniform on $[0^\circ,\,90^\circ]$, with $P$, $e$, and
$\omega_{\mathrm{peri}}$ fixed. No pipeline receives the injected
geometry or another pipeline's fit. Convergence statistics computed by
this project are reported for every comparator as information and never
block a result; failures are reported beside the successes. Every
method closes all reporting bins before the truth is opened.

\subsubsection{Reporting support and evaluation metrics}

% Question: On what wavelength support are methods compared, and with which
% metrics?
% Answer: 147 whole bins of a common R=100 grid inside the validated
% wavelength calibration, scored by RMSE, mean bias and shape RMS after
% closure of every bin.
% Evidence: products/w17_frozen_postclosure_score_20260923_r1_prepared/source/
% score_w17_frozen147.py and original_score_reference.py, and config/
% CASE_TEMPLATE_UNBOUND.json (regions); products/control_B_scored_20260921_r1/
% COMPARISON.npz; AGENTS.md (8 September decision); reports/
% W17_FROZEN_POSTCLOSURE_SCORE_PREPARATION_20260923_R1.md.

Recovered and injected spectra are compared on a common $R=100$ grid of
150 bins from 0.63 to 2.81~$\mu\mathrm{m}$. NOVA's native $2\times160$
coordinates are mapped to the grid by a fixed overlap map that requires
complete coverage of every retained bin. Bins 147 and 148 are excluded
because their detector support lies beyond the validated support of the
public wavelength calibration, where the solution is an extrapolation,
and bin 149 is excluded structurally. The retained support is bins 0 to
146, 147 bins, the last ending at $2.740018\,\mu\mathrm{m}$. This rule is
the same for every method, is applied only for reporting, and admits no
interpolation or gap filling: a method without an estimate in a retained
bin reports that bin as missing.

With $e_i=\widehat D_i-D_i^{\mathrm{inj}}$ on the retained bins, the
scores are
\begin{equation}
 \mathrm{RMSE}=
 \left(\frac{1}{N}\sum_i e_i^2\right)^{1/2},
 \label{eq:absolute-rmse}
\end{equation}
\begin{equation}
 \mathrm{RMS}_{\mathrm{shape}}=
 \left[\frac{1}{N}\sum_i(e_i-\bar e)^2\right]^{1/2},
 \label{eq:morphology-rmse}
\end{equation}
together with the mean bias $\bar e$ and the maximum absolute error, so
that $\mathrm{RMSE}^2=\bar e^{\,2}+\mathrm{RMS}_{\mathrm{shape}}^2$
separates a constant offset from an incorrect shape. No offset
correction is applied to a reported spectrum. The same scores are
reported on six contiguous ranges with boundaries at 1.0, 1.6, 1.8, 2.3,
and 2.7~$\mu\mathrm{m}$, and on two summary ranges, below
2.1~$\mu\mathrm{m}$ and 2.12 to 2.76~$\mu\mathrm{m}$. The injected
spectrum is opened only after a recovery's prediction on all 147 bins is
recorded; until then every calibration, modelling, and starting choice
rests on the data and the instrument references alone
(Appendix~\ref{app:benchmarking}). Each score is reported with the
standardised residuals, the conditional covariance, and the weight,
conditioning, and convergence diagnostics of the fit that produced it;
the validated uncertainty comes from the ensemble of
Section~\ref{sec:ensemble}.
